\chapter{Đào sâu: Khả năng quan sát --- trace xuyên Feign và Kafka, lấy mẫu, và những cảnh báo đáng giá}
\label{ch:dd-observability}

Chương~24 dạy bạn debug stack từ ngoài vào trong --- \code{describe}
trước \code{logs} --- và Chương~27 gọi tên bộ ba tiếp theo: metric,
trace, cảnh báo. Chương đào sâu này mở toang phần dây nối đã có sẵn
trong mọi service và hỏi nó thực sự làm gì. Câu trả lời ít hơn cấu hình
gợi ý: dependency tracing nằm trên mọi classpath, mẫu log in ra trace
id, cụm lấy mẫu ở mức không, và endpoint Prometheus mà Chương~27 trông
cậy thì chưa tồn tại. Đọc lý do vì sao là con đường nhanh nhất để hiểu
truy vết phân tán hoạt động thế nào.

Ba câu hỏi sắp xếp chương này:

\begin{enumerate}
  \item \textbf{Một trace bắt đầu ở đâu, và sống sót thế nào} qua các
  chặng trình duyệt $\rightarrow$ gateway $\rightarrow$ Feign
  $\rightarrow$ Kafka $\rightarrow$ consumer --- và ở đâu trong dự án
  này nó lặng lẽ chết?
  \item \textbf{Lấy mẫu tốn gì} ở 1.0, ở 0.0, và ở những con số hệ
  thống production thực sự dùng?
  \item \textbf{Tín hiệu nào xứng đáng một cảnh báo}, và phải thêm gì
  trước khi bất kỳ cảnh báo nào nổ được?
\end{enumerate}

\section{Vì sao chọn nó, và nó giải quyết gì}

Log trả lời ``tiến trình này đã nói gì?'' Metric trả lời ``hệ thống
đang ra sao, tính gộp, theo thời gian?'' Trace trả lời câu hỏi hai thứ
kia không thể: ``chuyện gì đã xảy ra với \emph{đúng một request này}
khi nó đi qua năm service và một message broker?'' Trong monolith,
một stack trace chính là một trace. Trong dự án này, một cú click
``tạo bài tập'' chạm tới gateway, course-service, auth-service (qua
Feign), Kafka, và notification-service; không file log đơn lẻ nào
thấy hết.

Dự án chọn Micrometer Tracing với cầu nối Brave, báo cáo về Zipkin.
\code{pom.xml} của mỗi service mang cùng một cặp:

\begin{lstlisting}[language=yaml]
<dependency>
  <groupId>io.micrometer</groupId>
  <artifactId>micrometer-tracing-bridge-brave</artifactId>   (*@\co{1}@*)
</dependency>
<dependency>
  <groupId>io.zipkin.reporter2</groupId>
  <artifactId>zipkin-reporter-brave</artifactId>             (*@\co{2}@*)
</dependency>
\end{lstlisting}

\begin{description}
  \item[\co{1}] Micrometer Tracing là \emph{facade} mà Spring Boot~3
  instrument lên; Brave (tracer của dự án Zipkin) là một trong hai
  hiện thực, cái còn lại là OpenTelemetry. Cầu nối là thứ biến các
  \emph{observation} của Spring --- Micrometer Observation API mà mọi
  starter Boot~3 phát ra --- thành span.
  \item[\co{2}] Reporter gửi các span đã hoàn tất tới Zipkin qua HTTP.
  Đổi đúng một artifact này lấy \code{opentelemetry-exporter-otlp} là
  cùng các span đó trỏ sang Tempo, Jaeger, hay một vendor.
\end{description}

\begin{decision}[Brave và Zipkin, không phải OpenTelemetry]
Brave cộng Zipkin là con đường ngắn nhất từ số không tới một màn hình
trace hoạt động: một container, không collector, không thương lượng
giao thức, và auto-configuration của Spring Boot biết cả hai rất rõ.
Đó là lựa chọn đúng để học cơ chế và cho một homelab. Điều sẽ lật nó:
một ngôn ngữ thứ hai trong stack (SDK OpenTelemetry phủ tất cả với một
giao thức đường truyền), một backend vendor (Datadog, Honeycomb,
Grafana Cloud đều nói OTLP bản địa), hoặc nhu cầu lấy mẫu theo đuôi
(tail-based), thứ đòi một collector đứng giữa ứng dụng và kho. Việc
chuyển đổi là cơ học --- đổi hai dependency, đổi thuộc tính endpoint
--- vì code chưa bao giờ chạm trực tiếp vào Brave. Người phỏng vấn hỏi
``Zipkin hay Jaeger hay Tempo?'' và mong đợi đúng câu trả lời này:
backend thay được, chuẩn instrumentation mới là quyết định tồn tại lâu.
\end{decision}

\section{Một trace đi qua dự án này như thế nào}

\subsection{Cấu hình, trong mọi service}

Năm tài liệu service trong config-server đều kết thúc giống nhau:

\begin{lstlisting}[language=yaml]
management:
  tracing:
    sampling:
      probability: 1.0                                   (*@\co{1}@*)
  zipkin:
    tracing:
      endpoint: http://localhost:9411/api/v2/spans       (*@\co{2}@*)
  endpoints:
    web:
      exposure:
        include: health,info,metrics,prometheus          (*@\co{3}@*)
logging:
  pattern:
    level: "%5p [${spring.application.name:},%X{traceId:-},%X{spanId:-}]"  (*@\co{4}@*)
\end{lstlisting}

\begin{description}
  \item[\co{1}] Mọi request đều được trace. Ổn trên laptop; xem
  Mục~\ref{sec:sampling} để biết nó tốn gì ở nơi khác. Trong cụm, mọi
  Deployment ghi đè thành \code{0.0} với chú thích ``Zipkin chưa có
  trong cụm''.
  \item[\co{2}] Compose ghi đè host thành \code{zipkin:9411}; reporter
  POST các lô JSON tới đây.
  \item[\co{3}] Exposure là một \emph{danh sách tên}. Nêu tên
  \code{prometheus} không tạo ra endpoint; endpoint chỉ tồn tại khi
  \code{micrometer-registry-prometheus} nằm trên classpath, và không
  \code{pom.xml} nào trong repository có nó. Hôm nay
  \code{/actuator/prometheus} trả 404 trên mọi service --- kế hoạch ở
  Chương~27 giả định một endpoint vẫn còn phải thêm vào.
  \item[\co{4}] Mẫu log đọc hai khóa MDC mà cầu nối tracing điền vào.
  Khi không có span đang hoạt động, giá trị dự phòng \code{-} in ra
  rỗng, và chính điều đó là dấu hiệu chẩn đoán: trace id rỗng trong
  một dòng log nghĩa là code chạy ngoài mọi ngữ cảnh được trace.
\end{description}

\subsection{Từng chặng một}

Theo dõi một \code{POST /api/courses/7/assignments} từ trình duyệt.

\textbf{Trình duyệt $\rightarrow$ gateway.} Trình duyệt không gửi
header trace nào; \code{axios} không biết gì về tracing. Gateway là ứng
dụng WebFlux và Spring Cloud Gateway được instrument qua Observation
API, nên observation phía server của request đi vào \emph{khởi tạo}
một trace mới: một trace id ngẫu nhiên 128 bit và một span gốc. Khi
gateway chuyển tiếp tới course-service, \code{WebClient} reactive nó
dùng cũng được instrument, và propagator ghi ngữ cảnh vào request đi
ra. Mặc định của Spring Boot~3 quan trọng ở đây: Brave được cấu hình
để \emph{phát} W3C Trace Context
(\code{traceparent: 00-<traceId>-<spanId>-01}) và \emph{đọc} cả W3C
\emph{lẫn} hai định dạng B3. Vậy dự án này nói W3C trên đường truyền
và vẫn hiểu một bên gọi B3 cũ hơn.

\textbf{Gateway $\rightarrow$ course-service.} Course-service là
servlet MVC. Observation phía server của nó đọc \code{traceparent},
tiếp tục cùng trace, và mở một span con. Mọi thứ log bên trong
controller giờ mang trace id đó --- mẫu ở \co{4} phía trên in nó ra.

\textbf{course-service $\rightarrow$ auth-service qua Feign.}
\code{AuthServiceClient} là một client Spring Cloud OpenFeign. OpenFeign
kèm sẵn một capability Micrometer, bật mặc định, bọc mỗi lời gọi trong
một observation; propagator của Brave tiêm \code{traceparent} vào
request Feign. Auth-service tiếp tục trace. Trong Zipkin, lời gọi Feign
hiện ra như một span client trong course-service lồng dưới span
controller, cùng một span server khớp trong auth-service. Nếu circuit
breaker của Resilience4j ngắt lời gọi (Chương~6), không span client
nào được tạo --- fallback đã chạy, đường truyền không hề bị chạm, và
trace hiện ra khoảng trống.

\textbf{course-service $\rightarrow$ Kafka.} Ở đây trace chết. Đọc
producer của dự án:

\begin{lstlisting}[language=Java]
@Bean
public ProducerFactory<String, Object> producerFactory() {
    Map<String, Object> configProps = new HashMap<>();
    configProps.put(BOOTSTRAP_SERVERS_CONFIG, bootstrapServers);
    configProps.put(KEY_SERIALIZER_CLASS_CONFIG, StringSerializer.class);
    configProps.put(VALUE_SERIALIZER_CLASS_CONFIG, JsonSerializer.class);
    return new DefaultKafkaProducerFactory<>(configProps);
}

@Bean
public KafkaTemplate<String, Object> kafkaTemplate() {
    return new KafkaTemplate<>(producerFactory());       (*@\co{1}@*)
}
\end{lstlisting}

\begin{description}
  \item[\co{1}] Một \code{KafkaTemplate} dựng tay. Spring Kafka
  \emph{có thể} lan truyền ngữ cảnh trace --- nó ghi \code{traceparent}
  vào header của record và mở một span producer --- nhưng chỉ khi
  \code{setObservationEnabled(true)} được gọi trên template, hoặc
  \code{spring.kafka.template.observation-enabled=true} được đặt cho
  bean tự cấu hình. Bean này không làm cả hai, và thuộc tính vắng mặt
  trong mọi tài liệu cấu hình. Record rời đi không có header. Phía tiêu
  thụ, \code{@KafkaListener} trong notification-service sẽ cần
  \code{spring.kafka.listener.observation-enabled=true} để đọc một
  header. Cả hai nửa đều tắt. E-mail sinh ra từ cú click, với Zipkin,
  là một sự kiện không liên quan.
\end{description}

\code{UserRegisteredPublisher} của auth-service dùng template tự cấu
hình và có cùng khoảng trống vì cùng lý do: thứ thiếu là thuộc tính,
không phải bean.

\textbf{Ranh giới scheduler.} \code{CourseLifecycleScheduler} chạy
\code{@Scheduled} lúc 01:00. Không gì instrument một trigger cron; dòng
log ``Auto-archived 3 course(s)'' của nó in ra với trace id và span id
rỗng. Điều tương tự đúng với mọi phương thức \code{@Async} trừ khi
executor được bọc bằng một decorator lan truyền ngữ cảnh. Trace là
thuộc tính của một \emph{request}; công việc không do request nào khởi
đầu thì không có trace trừ khi bạn trao cho nó một cái.

\subsection{Đối chiếu log}

Với mẫu ở \co{4}, một lệnh grep xuyên log của cả năm service theo một
trace id ghép lại câu chuyện Zipkin vẽ --- hữu ích đúng lúc Zipkin
chết hoặc bị lấy mẫu bỏ qua:

\begin{lstlisting}[language=bash]
$ docker compose logs --no-color api-gateway course-service \
    auth-service | grep 6f2c1a9e4b3d8e10
api-gateway     | INFO [api-gateway,6f2c1a9e4b3d8e10,9a1b...] ...
course-service  | INFO [course-service,6f2c1a9e4b3d8e10,c04d...] ...
auth-service    | INFO [auth-service,6f2c1a9e4b3d8e10,e77f...] ...
\end{lstlisting}

Để ý thứ vắng mặt: không có dòng \code{notification-service}, vì không
trace id nào vượt qua Kafka. Lệnh grep cho thấy lỗ hổng đúng như Zipkin
sẽ cho thấy.

\section{Lấy mẫu}
\label{sec:sampling}

\code{probability: 1.0} nghĩa là mọi trace được ghi và mọi span được
gửi. Chi phí tuyến tính theo lưu lượng: mỗi span là vài trăm byte JSON,
được gom lô và POST bởi một thread nền, rồi được Zipkin đánh index và
lưu. Ở lưu lượng của dự án, chẳng đáng gì. Ở 1.000 request mỗi giây
với sáu span mỗi request, đó là 6.000 span mỗi giây, cỡ 2\,MB/s dữ liệu
vào và hàng chục gigabyte lưu trữ mỗi ngày --- phần lớn ghi lại những
request thành công trong 20\,ms mà không ai sẽ nhìn tới.

Lấy mẫu \emph{theo đầu} (head-based) --- thứ \code{probability} cấu
hình --- quyết định tại span gốc có giữ trace hay không, và lan truyền
quyết định trong cờ của \code{traceparent} để mọi service phía sau đồng
ý. Nó rẻ và nhất quán, và nó mù: ở 0.01 nó giữ 1\% trace chọn ngẫu
nhiên, nên request duy nhất thất bại lúc 3 giờ sáng được giữ với xác
suất 0.01. Lấy mẫu \emph{theo đuôi} (tail-based) hoãn quyết định tới
khi biết toàn bộ trace và giữ mọi trace có lỗi hoặc vượt ngân sách độ
trễ, cộng một mẫu ngẫu nhiên phần còn lại. Nó cần một thành phần đệm
trọn trace trước kho --- một OpenTelemetry Collector --- và đó là lý do
nó ngoài tầm với của dây nối trực tiếp Brave-tới-Zipkin.

\begin{decision}[tỷ lệ lấy mẫu]
Laptop và Compose: \code{1.0}, vì bạn đang debug và lưu lượng chính là
bạn. Cụm hôm nay: \code{0.0}, thành thật về việc chưa có pod Zipkin ---
lấy mẫu ở mức không cũng tắt cả chi phí tạo span, không chỉ chi phí
gửi. Cụm có Zipkin: bắt đầu ở \code{0.1}; hạ xuống \code{0.01} khi kho
lớn nhanh hơn tốc độ bạn đọc nó. Câu trả lời senior cho ``tỷ lệ lấy mẫu
bao nhiêu?'' là ``head-based 1--10\% làm nền, tail-based cho lỗi và
request chậm khi đã có collector'' --- và câu tiếp theo luôn là ``mà
trace lỗi mới là thứ bạn thực sự muốn, nên hãy tính ngân sách cho
collector''.
\end{decision}

\section{Metric: cái gì có và cái gì chưa}

\code{/actuator/metrics} của Actuator đã tồn tại trên mọi service và
liệt kê vài trăm meter: bộ nhớ JVM và GC, request HTTP server gắn tag
\code{uri}, \code{status} và \code{method}, metric client Kafka, trạng
thái pool HikariCP, và trạng thái circuit breaker của Resilience4j
trong course-service. Nó trả lời từng meter một, theo tên, qua HTTP ---
một công cụ debug, không phải công cụ giám sát. Giám sát nghĩa là một
hệ thống \emph{kéo} tất cả chúng mỗi 15 giây, lưu thành chuỗi thời
gian, và đánh giá quy tắc. Hệ thống đó là Prometheus, và endpoint của
nó đang thiếu (xem \co{3} ở trên).

Hai bộ từ vựng mang vào phỏng vấn. \textbf{RED} dành cho service theo
request: Rate, Errors, Duration --- ba thứ mà timer
\code{http.server.requests} đã cung cấp. \textbf{USE} dành cho tài
nguyên: Utilization, Saturation, Errors --- CPU, heap, pool Hikari,
consumer lag của Kafka. Một dashboard mỗi service nên hiện RED ở trên
và USE bên dưới; cảnh báo nên nổ trên một triệu chứng RED mà người
dùng cảm nhận được, và các panel USE giải thích vì sao.

\section{Ưu và nhược}

\begin{center}
\begin{tabular}{@{}p{0.47\linewidth}p{0.47\linewidth}@{}}
\toprule
\textbf{Điểm mạnh} & \textbf{Điểm yếu, hôm nay} \\
\midrule
Không đổi code: mọi chặng được các starter instrument, và mẫu log chỉ
một dòng. &
Lan truyền qua Kafka tắt ở cả producer lẫn consumer; con đường thú vị
nhất của hệ thống là con đường vô hình. \\
Facade trung lập vendor: Brave và Zipkin đổi được sang OTLP mà không
chạm mã nguồn. &
Không registry Prometheus trên classpath nào; \code{prometheus} trong
danh sách exposure chỉ là ước vọng. \\
\code{traceparent} W3C trên đường truyền, chấp nhận B3 --- tương tác
được với mọi thứ hiện đại. &
Chỉ lấy mẫu head-based; một lỗi ở mức 1\% được giữ 1\% số lần. \\
Đối chiếu log hoạt động cả khi không có Zipkin. &
Zipkin chạy \code{:latest} không volume: trace bốc hơi mỗi lần
\code{compose down}. \\
\bottomrule
\end{tabular}
\end{center}

\section{Cái gì gãy}

\begin{pitfall}[trace kết thúc ở Kafka]
Triệu chứng: trong Zipkin, trace \code{POST /assignments} kết thúc ở
span controller; thông báo theo sau là một trace một-span riêng biệt,
hoặc không có. Nguyên nhân: cả \code{spring.kafka.template.\allowbreak
observation-enabled} lẫn \code{spring.kafka.listener.\allowbreak
observation-enabled} đều không được đặt, nên không header
\code{traceparent} nào được ghi vào hay đọc từ record. Cách sửa: cả
hai thuộc tính, trong tài liệu cấu hình của cả hai service
(course-service và auth-service phát; notification-service tiêu thụ)
--- và template dựng tay trong \code{KafkaProducerConfig} phải gọi
\code{setObservationEnabled(true)} vì thuộc tính chỉ cấu hình bean tự
cấu hình. Listing ở Mục~\ref{sec:enhance} làm cả ba việc.
\end{pitfall}

\begin{pitfall}[tag có số lượng giá trị không giới hạn]
Triệu chứng: heap của một service leo dần suốt nhiều ngày và
\code{/actuator/metrics} mất vài giây để trả lời; cuối cùng
\code{OOMKilled}. Nguyên nhân: một tag có tập giá trị không bị chặn.
Kinh điển là một handler ánh xạ \code{/courses/\{id\}} bị ai đó viết
lại để dựng đường dẫn bằng tay --- tag \code{uri} khi đó mang
\code{/courses/42}, \code{/courses/43}, \ldots, mỗi cái một chuỗi thời
gian, mãi mãi. Chuyện tương tự xảy ra khi span được gắn tag bằng user
id hay e-mail. Cách sửa: gắn tag bằng \emph{mẫu} (\code{/courses/\{id\}}),
không bao giờ bằng giá trị; đặt định danh vào log annotation của span
hay dòng log, những thứ không được đánh index theo giá trị. Tag
\code{uri} mặc định của Spring là mẫu đường dẫn chính để tránh điều
này.
\end{pitfall}

\begin{pitfall}[một 401 không có trace]
Triệu chứng: người dùng báo ``tôi bị đăng xuất'' và không có trace nào
cho request của họ. Nguyên nhân: \code{JwtAuthenticationFilter} của
gateway từ chối ở thứ tự \code{-1}, trước khi định tuyến; observation
phía server vẫn ghi request, nhưng với lấy mẫu ở 0.0 trong cụm thì
không gì được xuất đi đâu cả. Cách sửa không phải tracing: đây là câu
hỏi \emph{metric}. Đếm \code{http.server.requests} với
\code{status=401} tại gateway và cảnh báo khi tăng vọt --- một lần
deploy sai \code{JWT\_SECRET} hiện ra như một vách đá trên bộ đếm đó
trong vòng 15 giây.
\end{pitfall}

\section{Chuyện gì gãy lúc 3 giờ sáng}

Zipkin chết lúc 3 giờ sáng; các service vẫn chạy.

\begin{itemize}
  \item \textbf{Độ trễ request: không đổi.} Reporter là bất đồng bộ.
  Span hoàn tất đi vào một hàng đợi trong bộ nhớ; một thread nền rút nó
  ra theo lô qua HTTP. Thread request không bao giờ chờ Zipkin. Đây là
  tính chất cần nói đầu tiên trong phỏng vấn, vì nỗi sợ ``tracing làm
  chậm service của tôi'' là phản đối phổ biến.
  \item \textbf{Bộ nhớ: bị chặn.} Hàng đợi có dung lượng cố định
  (\code{queuedMaxSpans}, mặc định 10.000). Khi đầy, span mới bị
  \emph{bỏ}, được đếm, và một cảnh báo được log --- reporter không
  phình heap để giữ chúng. Bạn mất tầm nhìn, không mất service.
  \item \textbf{Log: vẫn đối chiếu được.} Trace id được sinh cục bộ;
  mẫu MDC vẫn in chúng. Lệnh grep ở mục đối chiếu hoạt động đúng như
  trước. Đây là lý lẽ cho việc luôn có mẫu log, bất kể backend trace.
  \item \textbf{Hồi phục: tự động.} Reporter thử lại mỗi lô ở lần rút
  kế tiếp; khi Zipkin trở lại, span mới chảy tiếp. Các span bị bỏ trong
  lúc sự cố mất luôn --- không có replay.
\end{itemize}

Giờ theo chiều ngược lại: \emph{lưu lượng} tăng vọt lúc 3 giờ sáng với
lấy mẫu ở 1.0. Hàng đợi đầy, span bị bỏ, và chính dòng log cảnh báo
trở thành tiếng ồn ở hàng trăm dòng mỗi giây. Đó là hình hài cụ thể của
đánh đổi lấy mẫu --- và lý do đặt tỷ lệ trước cơn tăng vọt, không phải
trong lúc đó.

\section{Nâng cấp giải pháp}
\label{sec:enhance}

Bốn bổ sung biến dây nối thành một hệ thống giám sát. Mỗi cái đều trọn
vẹn; thêm theo thứ tự.

\textbf{1. Endpoint Prometheus.} Một dependency mỗi service (BOM Boot
của parent quản lý phiên bản):

\begin{lstlisting}[language=yaml]
<dependency>
  <groupId>io.micrometer</groupId>
  <artifactId>micrometer-registry-prometheus</artifactId>
</dependency>
\end{lstlisting}

Có nó trên classpath, danh sách exposure hiện có làm
\code{/actuator/prometheus} thành thật. Thêm các công tắc tracing đang
thiếu, và một tag chung để mọi chuỗi nói rõ nó đến từ service nào:

\begin{lstlisting}[language=yaml]
management:
  prometheus:
    metrics:
      export:
        enabled: true
  metrics:
    tags:
      application: ${spring.application.name}   # every series, one tag
  tracing:
    sampling:
      probability: ${TRACING_SAMPLE:0.1}          # env-overridable
spring:
  kafka:
    template:
      observation-enabled: true    # producer span + traceparent header
    listener:
      observation-enabled: true    # consumer span, continues the trace
\end{lstlisting}

Và template dựng tay trong course-service, thứ bỏ qua thuộc tính:

\begin{lstlisting}[language=Java]
@Bean
public KafkaTemplate<String, Object> kafkaTemplate() {
    KafkaTemplate<String, Object> template =
            new KafkaTemplate<>(producerFactory());
    // Properties configure only the auto-configured template; a
    // hand-built one must opt in explicitly or traces end here.
    template.setObservationEnabled(true);
    return template;
}
\end{lstlisting}

\textbf{2. Scrape trong cụm.} Với \code{kube-prometheus-stack} đã cài
(Chương~27), một \code{ServiceMonitor} bảo operator của nó scrape
Service nào. Một object phủ cả năm, theo label:

\begin{lstlisting}[language=yaml]
apiVersion: monitoring.coreos.com/v1
kind: ServiceMonitor
metadata:
  name: ts-services
  namespace: ts
  labels:
    release: kube-prometheus-stack   # the selector the chart installs
spec:
  selector:
    matchLabels:
      monitoring: actuator           # add this label to each Service
  endpoints:
    - port: http                     # the Service port *name*
      path: /actuator/prometheus
      interval: 15s
\end{lstlisting}

Service của auth-service và course-service phải cho phép
\code{/actuator/prometheus} trong SecurityConfig của chúng để việc này
chạy --- cùng phép cho phép mà Chương~17 cần cho probe. Hạn chế nó
trong mạng pod bằng NetworkPolicy thay vì để mở ra Ingress.

\textbf{3. Log có cấu trúc.} Grep xuyên container hoạt động trên
laptop; trong cụm, Loki hay Elasticsearch parse JSON đáng tin hơn nhiều
so với một mẫu ngoặc vuông. Bộ mã hóa JSON của Logback là một dependency
cộng một \code{logback-spring.xml}:

\begin{lstlisting}[language=yaml]
<dependency>
  <groupId>net.logstash.logback</groupId>
  <artifactId>logstash-logback-encoder</artifactId>
  <version>8.0</version>
</dependency>
\end{lstlisting}

\begin{lstlisting}
<configuration>
  <springProperty name="app" source="spring.application.name"/>
  <appender name="JSON"
            class="ch.qos.logback.core.ConsoleAppender">
    <encoder class="net.logstash.logback.encoder.LogstashEncoder">
      <customFields>{"application":"${app}"}</customFields>
      <!-- traceId and spanId come from the MDC automatically -->
      <includeMdcKeyName>traceId</includeMdcKeyName>
      <includeMdcKeyName>spanId</includeMdcKeyName>
    </encoder>
  </appender>
  <springProfile name="!local">
    <root level="INFO"><appender-ref ref="JSON"/></root>
  </springProfile>
  <springProfile name="local">
    <include resource="org/springframework/boot/logging/logback/base.xml"/>
  </springProfile>
</configuration>
\end{lstlisting}

Việc tách theo profile giữ mẫu đọc-được-bằng-mắt trên laptop và phát
một object JSON mỗi dòng ở mọi nơi khác. Mỗi dòng giờ mang
\code{traceId} như một trường; truy vấn log theo trace id là một lượt
tra index, không phải quét.

\textbf{4. Hai cảnh báo gắn với điều người dùng cảm nhận.} Quy tắc
Prometheus, trong kind \code{PrometheusRule} của operator:

\begin{lstlisting}[language=yaml]
apiVersion: monitoring.coreos.com/v1
kind: PrometheusRule
metadata:
  name: ts-slo
  namespace: ts
  labels: { release: kube-prometheus-stack }
spec:
  groups:
    - name: ts.slo
      rules:
        - alert: GatewayErrorBudgetBurn
          expr: |
            sum(rate(http_server_requests_seconds_count{
                  application="api-gateway", status=~"5.."}[5m]))
            /
            sum(rate(http_server_requests_seconds_count{
                  application="api-gateway"}[5m])) > 0.02
          for: 10m
          labels: { severity: page }
          annotations:
            summary: "Gateway 5xx ratio above 2% for 10 minutes"
        - alert: NotificationConsumerLag
          expr: |
            sum by (topic) (
              kafka_consumer_fetch_manager_records_lag_max{
                application="notification-service"}) > 500
          for: 15m
          labels: { severity: ticket }
          annotations:
            summary: "notification-group is {{ $value }} records behind"
\end{lstlisting}

Cái đầu là cảnh báo \emph{triệu chứng} trên tỷ lệ ``errors'' của RED
tại nơi duy nhất mọi request đi qua; 2\% trong mười phút là ngân sách
khởi điểm, không phải luật. Cái thứ hai là câu trả lời của Chương~27
cho ``metric Kafka nào trước'': lag tăng suốt 15 phút nghĩa là mail
không đi ra, bất kể sức khỏe của pod nói gì. Để ý tên meter: Spring
Kafka đăng ký \code{records-lag-max} của chính client theo cách đặt tên
của Micrometer, và registry Prometheus hiển thị nó bằng gạch dưới.

\begin{handson}[chứng minh trace vượt qua Kafka]
Sau khi thêm hai thuộc tính \code{observation-enabled} và dòng
template, khởi động lại course-service và notification-service dưới
Compose, tạo một bài tập qua UI, rồi hỏi Zipkin về trace theo id in
trong log của course-service:
\begin{lstlisting}[language=bash]
$ curl -s "http://localhost:9411/api/v2/trace/6f2c1a9e4b3d8e10" \
    | jq -r '.[] | "\(.localEndpoint.serviceName)  \(.name)"'
api-gateway            http post
course-service         http post /courses/{id}/assignments
course-service         http get                   # Feign to auth
auth-service           http get /users/{id}
course-service         ts.assignment.created send
notification-service   ts.assignment.created process
\end{lstlisting}
Trước thay đổi, hai dòng cuối không tồn tại. Span consumer có cha là
span producer; e-mail giờ quy được về cú click đã gây ra nó.
\end{handson}

\section{Ngoài dự án}

Các stack production hội tụ về cùng một hình dạng với những cái tên
khác nhau. Instrumentation: SDK OpenTelemetry hoặc Java agent, thứ tự
instrument mà không cần cả starter. Thu thập: một OpenTelemetry
Collector mỗi node hay mỗi cụm, nơi lấy mẫu tail-based, che dữ liệu,
và phân phối tới nhiều backend cư trú. Lưu trữ: Tempo cho trace, Loki
cho log, Prometheus hay Mimir cho metric --- bộ ba Grafana --- hoặc
Jaeger, Elasticsearch và Thanos; hoặc một vendor nhận OTLP và tính tiền
theo dung lượng. Mạch phỏng vấn chạy xuyên tất cả: ``làm sao tìm một
request chậm xuyên các service?'' được trả lời bằng trace id trong log
và một kho trace; ``làm sao biết trước khi người dùng báo?'' bằng cảnh
báo RED ở biên; ``tracing tốn gì?'' bằng thảo luận lấy mẫu ở trên ---
và phiên bản senior của mọi câu trả lời gọi tên thành phần ở lại
(instrumentation và trace id trong mọi dòng log) và những thành phần
thay được (mọi thứ phía sau exporter).

\section{Câu hỏi từ phỏng vấn thật}

\subsection*{Q: Độ trễ p99 tăng gấp đôi đêm qua. Hãy dẫn tôi qua cách tìm nguyên nhân.}

Bắt đầu từ metric, không phải log: histogram
\code{http\_server\_requests} tại gateway, tách theo \code{uri}, cho
bạn biết trong một truy vấn liệu \emph{mọi} route đều chậm (hạ tầng
--- node, mạng, GC) hay chỉ một route (một nhánh code hoặc một
downstream). Rồi lấy một trace chậm từ khung giờ đó --- với lấy mẫu
head-based ở 1\% sẽ có vài trăm --- và đọc thác span: span chậm hoặc là
một span con (một lời gọi downstream, chẳng hạn chặng Feign tới
auth-service mất 800\,ms thay vì 8) hoặc là khoảng trống giữa các span
con (chính service đó, thường là một lời gọi cơ sở dữ liệu hay một
lock, hiện ra như thời gian trong span cha không được span con nào phủ).
Nếu là khoảng trống, trace id đưa bạn tới log của service đó đúng phút
ấy; layout JSON làm việc đó thành một lượt tra index. Chi tiết senior
là thứ tự --- tổng hợp, rồi một mẫu đại diện, rồi log của nó --- và gọi
tên thứ bạn ước có: một liên kết exemplar từ bucket histogram thẳng tới
một trace, thứ các stack gốc OpenTelemetry cung cấp còn Zipkin thì
không.

\subsection*{Q: Bạn không thể trace 100\%. Bạn lấy mẫu gì, và không bao giờ bỏ gì?}

Lấy mẫu head-based 1--10\% cho trạng thái ổn định, quyết định tại
gateway và lan truyền trong cờ của \code{traceparent} để các service
phía sau không bao giờ bất đồng. Không bao giờ bỏ lỗi và không bao giờ
bỏ phần đuôi chậm --- điều lấy mẫu head-based không thể hứa, vì span
gốc được lấy mẫu trước khi ai biết request sẽ thất bại. Nên câu trả lời
thật có hai tầng: một collector lấy mẫu tail-based giữ mọi trace có 5xx
hoặc vượt ngân sách độ trễ, cộng một nền ngẫu nhiên 1\% cho hình dáng
lưu lượng bình thường. Nói thẳng về chi phí: lấy mẫu tail-based đệm
trọn trace trong bộ nhớ collector suốt cửa sổ quyết định, thường
10--30\,s, nên collector được định cỡ theo lưu lượng, không theo thứ
bạn giữ. Trong dự án này đó là bước từ \code{zipkin-reporter-brave}
sang một exporter OTLP cộng một collector; code ứng dụng không đổi.

\subsection*{Q: Một request biến mất khỏi trace ở Kafka. Vì sao, và sửa thế nào?}

Vì ngữ cảnh trace chỉ vượt qua Kafka dưới dạng header của record, và
không gì ghi chúng trừ khi template được bảo. Trong Spring Kafka hai
thuộc tính làm việc đó, cả hai mặc định false trong Boot~3: phía
producer, \code{template.observation-enabled}; phía consumer,
\code{listener.observation-enabled} (cả hai dưới \code{spring.kafka}).
Điểm tinh tế đáng nói to: thuộc tính chỉ cấu hình template tự cấu hình,
nên một bean \code{KafkaTemplate} dựng tay --- đúng thứ
\code{KafkaProducerConfig} của course-service làm --- cần
\code{setObservationEnabled(true)} trong code. Sau khi sửa, span
consumer có cha là span producer và toàn bộ đường e-mail treo dưới
request HTTP gốc. Câu tiếp theo người phỏng vấn muốn: span consumer là
\emph{con}, nhưng nó có thể chạy vài giây hay vài giờ sau khi request
kết thúc, nên màn hình trace sắp theo thời gian sẽ hiện một khoảng
trống dài --- khoảng trống đó là độ trễ hàng đợi, và bản thân nó là
metric cần theo dõi.

\subsection*{Q: Bạn sẽ đánh thức người trực vì cái gì với stack này?}

Triệu chứng người dùng cảm nhận, diễn đạt thành SLO, không bao giờ là
tài nguyên thô. Hai cảnh báo gọi người: tỷ lệ 5xx của gateway vượt ngân
sách lỗi suốt mười phút --- điểm duy nhất mọi request đi qua --- và độ
trễ p99 tại gateway vượt mục tiêu trong cùng cửa sổ. Một cảnh báo mức
ticket: consumer lag của notification tăng suốt mười lăm phút, vì đó là
``mail không đi ra'' ở dạng duy nhất hệ thống diễn đạt được. Mọi thứ
khác --- CPU, heap, pod khởi động lại, pool Hikari cạn --- là panel
dashboard giải thích một cảnh báo, không phải bản thân cảnh báo; cảnh
báo CPU 80\% đánh thức người vì một hệ thống đang hoạt động. Khung
senior là ngân sách lỗi: ở mục tiêu sẵn sàng 99,9\% bạn được đốt 43
phút lỗi mỗi tháng, và cảnh báo nổ trên \emph{tốc độ đốt}, đó là lý do
tỷ lệ được đo trong một cửa sổ chứ không trên một request thất bại
đơn lẻ.

\subsection*{Q: Request của một người dùng chạm năm service. Làm sao tìm năm dòng log đó?}

Bằng trace id, thứ mọi service đều in vì mẫu log có
\code{\%X\{traceId\}} và cầu nối tracing điền MDC ở mọi chặng. Trên
laptop đó là grep xuyên log container; trong cụm là truy vấn trên một
kho log đánh index trường đó, và đó là lý do log được gửi dưới dạng
JSON thay vì parse bằng regex. Hai chi tiết tách một câu trả lời
senior: thứ nhất, trace id phải được \emph{trả về cho người dùng} ---
một header phản hồi \code{X-Trace-Id} từ gateway --- để một ticket hỗ
trợ trích dẫn được nó thay vì ``khoảng 3 giờ chiều''; thứ hai, bất kỳ
thứ gì chạy ngoài request (\code{@Scheduled}, \code{@Async}, một
consumer Kafka không lan truyền) log ra id rỗng, và bạn nên biết những
nhánh code nào của mình thuộc loại đó trước sự cố.

\subsection*{Q: Zipkin chết. Production có chậm đi không?}

Không, và bạn nên nói được vì sao từ thiết kế của reporter: span được
trao cho một hàng đợi trong bộ nhớ và một thread nền POST chúng theo
lô, nên không thread request nào chờ Zipkin; hàng đợi bị chặn ở 10.000
span và bỏ khi tràn, nên bộ nhớ có trần; các lô thất bại được log và
thử lại ở lần rút kế tiếp, nên hồi phục là tự động. Thứ bạn mất là span
trong cửa sổ sự cố --- không có replay. Cái bẫy trong câu hỏi là trường
hợp ngược lại: một exporter \emph{đồng bộ}, hay một hàng đợi không
giới hạn, sẽ biến sự cố giám sát thành sự cố ứng dụng, và đó là tính
chất cần kiểm tra trước khi đưa vào bất kỳ agent hay exporter nào. Và
cách giảm nhẹ cho các span bị mất cũng giống mọi thứ khác trong chương
này: trace id vẫn nằm trong mọi dòng log, nên log mang sự đối chiếu
ngay cả khi kho trace không thể.

\begin{quiz}
\begin{enumerate}
  \item Nêu chính xác thuộc tính hay lời gọi phương thức đang thiếu ở
  mỗi trong hai nửa phía Kafka của dự án, và giải thích vì sao chỉ đặt
  thuộc tính sẽ sửa được auth-service nhưng không sửa được
  course-service.
  \item Một request tới gateway mang \code{X-B3-TraceId} từ một proxy
  upstream cũ. Dự án tiếp tục trace đó hay khởi tạo trace mới, và mặc
  định nào của Boot~3 quyết định điều đó? Header nào rời gateway hướng
  về course-service?
  \item \code{/actuator/prometheus} nằm trong mọi danh sách exposure
  mà vẫn trả 404. Giải thích, và nêu tên đúng một artifact thay đổi câu
  trả lời.
  \item Zipkin chết một giờ dưới lưu lượng bình thường. Cho biết chuyện
  gì xảy ra với độ trễ request, với heap, và với các span sinh ra trong
  giờ đó --- và cái nào trong ba thứ sẽ đổi nếu reporter là đồng bộ.
  \item Ở \code{probability: 0.01}, tỷ lệ request \emph{thất bại} có
  trace là bao nhiêu, và thành phần kiến trúc nào cần có để đưa tỷ lệ
  đó lên 1.0 mà không trace mọi request thành công?
  \item Một đồng nghiệp thêm \code{span.tag("email", user.getEmail())}
  vào span đăng nhập. Dự đoán sự cố trong những tuần tiếp theo và nêu
  quy tắc ngăn nó.
\end{enumerate}
\end{quiz}
